% Options for packages loaded elsewhere
\PassOptionsToPackage{unicode}{hyperref}
\PassOptionsToPackage{hyphens}{url}
\documentclass[
]{article}
\usepackage{xcolor}
\usepackage[margin=1in]{geometry}
\usepackage{amsmath,amssymb}
\setcounter{secnumdepth}{-\maxdimen} % remove section numbering
\usepackage{iftex}
\ifPDFTeX
  \usepackage[T1]{fontenc}
  \usepackage[utf8]{inputenc}
  \usepackage{textcomp} % provide euro and other symbols
\else % if luatex or xetex
  \usepackage{unicode-math} % this also loads fontspec
  \defaultfontfeatures{Scale=MatchLowercase}
  \defaultfontfeatures[\rmfamily]{Ligatures=TeX,Scale=1}
\fi
\usepackage{lmodern}
\ifPDFTeX\else
  % xetex/luatex font selection
\fi
% Use upquote if available, for straight quotes in verbatim environments
\IfFileExists{upquote.sty}{\usepackage{upquote}}{}
\IfFileExists{microtype.sty}{% use microtype if available
  \usepackage[]{microtype}
  \UseMicrotypeSet[protrusion]{basicmath} % disable protrusion for tt fonts
}{}
\makeatletter
\@ifundefined{KOMAClassName}{% if non-KOMA class
  \IfFileExists{parskip.sty}{%
    \usepackage{parskip}
  }{% else
    \setlength{\parindent}{0pt}
    \setlength{\parskip}{6pt plus 2pt minus 1pt}}
}{% if KOMA class
  \KOMAoptions{parskip=half}}
\makeatother
\usepackage{color}
\usepackage{fancyvrb}
\newcommand{\VerbBar}{|}
\newcommand{\VERB}{\Verb[commandchars=\\\{\}]}
\DefineVerbatimEnvironment{Highlighting}{Verbatim}{commandchars=\\\{\}}
% Add ',fontsize=\small' for more characters per line
\usepackage{framed}
\definecolor{shadecolor}{RGB}{248,248,248}
\newenvironment{Shaded}{\begin{snugshade}}{\end{snugshade}}
\newcommand{\AlertTok}[1]{\textcolor[rgb]{0.94,0.16,0.16}{#1}}
\newcommand{\AnnotationTok}[1]{\textcolor[rgb]{0.56,0.35,0.01}{\textbf{\textit{#1}}}}
\newcommand{\AttributeTok}[1]{\textcolor[rgb]{0.13,0.29,0.53}{#1}}
\newcommand{\BaseNTok}[1]{\textcolor[rgb]{0.00,0.00,0.81}{#1}}
\newcommand{\BuiltInTok}[1]{#1}
\newcommand{\CharTok}[1]{\textcolor[rgb]{0.31,0.60,0.02}{#1}}
\newcommand{\CommentTok}[1]{\textcolor[rgb]{0.56,0.35,0.01}{\textit{#1}}}
\newcommand{\CommentVarTok}[1]{\textcolor[rgb]{0.56,0.35,0.01}{\textbf{\textit{#1}}}}
\newcommand{\ConstantTok}[1]{\textcolor[rgb]{0.56,0.35,0.01}{#1}}
\newcommand{\ControlFlowTok}[1]{\textcolor[rgb]{0.13,0.29,0.53}{\textbf{#1}}}
\newcommand{\DataTypeTok}[1]{\textcolor[rgb]{0.13,0.29,0.53}{#1}}
\newcommand{\DecValTok}[1]{\textcolor[rgb]{0.00,0.00,0.81}{#1}}
\newcommand{\DocumentationTok}[1]{\textcolor[rgb]{0.56,0.35,0.01}{\textbf{\textit{#1}}}}
\newcommand{\ErrorTok}[1]{\textcolor[rgb]{0.64,0.00,0.00}{\textbf{#1}}}
\newcommand{\ExtensionTok}[1]{#1}
\newcommand{\FloatTok}[1]{\textcolor[rgb]{0.00,0.00,0.81}{#1}}
\newcommand{\FunctionTok}[1]{\textcolor[rgb]{0.13,0.29,0.53}{\textbf{#1}}}
\newcommand{\ImportTok}[1]{#1}
\newcommand{\InformationTok}[1]{\textcolor[rgb]{0.56,0.35,0.01}{\textbf{\textit{#1}}}}
\newcommand{\KeywordTok}[1]{\textcolor[rgb]{0.13,0.29,0.53}{\textbf{#1}}}
\newcommand{\NormalTok}[1]{#1}
\newcommand{\OperatorTok}[1]{\textcolor[rgb]{0.81,0.36,0.00}{\textbf{#1}}}
\newcommand{\OtherTok}[1]{\textcolor[rgb]{0.56,0.35,0.01}{#1}}
\newcommand{\PreprocessorTok}[1]{\textcolor[rgb]{0.56,0.35,0.01}{\textit{#1}}}
\newcommand{\RegionMarkerTok}[1]{#1}
\newcommand{\SpecialCharTok}[1]{\textcolor[rgb]{0.81,0.36,0.00}{\textbf{#1}}}
\newcommand{\SpecialStringTok}[1]{\textcolor[rgb]{0.31,0.60,0.02}{#1}}
\newcommand{\StringTok}[1]{\textcolor[rgb]{0.31,0.60,0.02}{#1}}
\newcommand{\VariableTok}[1]{\textcolor[rgb]{0.00,0.00,0.00}{#1}}
\newcommand{\VerbatimStringTok}[1]{\textcolor[rgb]{0.31,0.60,0.02}{#1}}
\newcommand{\WarningTok}[1]{\textcolor[rgb]{0.56,0.35,0.01}{\textbf{\textit{#1}}}}
\usepackage{longtable,booktabs,array}
\usepackage{calc} % for calculating minipage widths
% Correct order of tables after \paragraph or \subparagraph
\usepackage{etoolbox}
\makeatletter
\patchcmd\longtable{\par}{\if@noskipsec\mbox{}\fi\par}{}{}
\makeatother
% Allow footnotes in longtable head/foot
\IfFileExists{footnotehyper.sty}{\usepackage{footnotehyper}}{\usepackage{footnote}}
\makesavenoteenv{longtable}
\usepackage{graphicx}
\makeatletter
\newsavebox\pandoc@box
\newcommand*\pandocbounded[1]{% scales image to fit in text height/width
  \sbox\pandoc@box{#1}%
  \Gscale@div\@tempa{\textheight}{\dimexpr\ht\pandoc@box+\dp\pandoc@box\relax}%
  \Gscale@div\@tempb{\linewidth}{\wd\pandoc@box}%
  \ifdim\@tempb\p@<\@tempa\p@\let\@tempa\@tempb\fi% select the smaller of both
  \ifdim\@tempa\p@<\p@\scalebox{\@tempa}{\usebox\pandoc@box}%
  \else\usebox{\pandoc@box}%
  \fi%
}
% Set default figure placement to htbp
\def\fps@figure{htbp}
\makeatother
\setlength{\emergencystretch}{3em} % prevent overfull lines
\providecommand{\tightlist}{%
  \setlength{\itemsep}{0pt}\setlength{\parskip}{0pt}}
\usepackage{bookmark}
\IfFileExists{xurl.sty}{\usepackage{xurl}}{} % add URL line breaks if available
\urlstyle{same}
\hypersetup{
  pdftitle={DATAX500},
  pdfauthor={1631819},
  hidelinks,
  pdfcreator={LaTeX via pandoc}}

\title{DATAX500}
\author{1631819}
\date{2026-03-22}

\begin{document}
\maketitle

\begin{longtable}[]{@{}
  >{\raggedright\arraybackslash}p{(\linewidth - 0\tabcolsep) * \real{1.0000}}@{}}
\toprule\noalign{}
\endhead
\bottomrule\noalign{}
\endlastfoot
\href{https://github.com/BrewBarred/DATAX500_A1}{Github Repository} \\
\href{https://elearn.waikato.ac.nz/mod/assign/view.php?id=2241094&action=editsubmission}{Moodle
Submission} \\
\end{longtable}

\section{Assignment 1}\label{assignment-1}

\begin{enumerate}
\def\labelenumi{\arabic{enumi}.}
\tightlist
\item
  Using the `\emph{heart.csv}' data-set, calculate the following
  statistical measures for age in \emph{R} \textbf{(6 marks)}
\end{enumerate}

\begin{itemize}
\item
  First, we load the data-set:

\begin{Shaded}
\begin{Highlighting}[]
\CommentTok{\# load automatically via git files}
\NormalTok{heart}\OtherTok{\textless{}{-}}\FunctionTok{read.csv}\NormalTok{(}\StringTok{"Heart.csv"}\NormalTok{)}


\DocumentationTok{\#\# Alternatively, load the dataset manually using:}
\CommentTok{\# heart\textless{}{-}read.csv(file.choose())}
\end{Highlighting}
\end{Shaded}
\end{itemize}

\begin{enumerate}
\def\labelenumi{\alph{enumi})}
\item
  Mean \textbf{(2 marks)}

\begin{Shaded}
\begin{Highlighting}[]
\CommentTok{\# assign age column to "age" var for easier processing}
\NormalTok{age}\OtherTok{\textless{}{-}}\NormalTok{heart}\SpecialCharTok{$}\NormalTok{age}

\CommentTok{\# calculate the mean age}
\NormalTok{mean\_age}\OtherTok{\textless{}{-}}\FunctionTok{mean}\NormalTok{(age)}
\NormalTok{mean\_age}
\end{Highlighting}
\end{Shaded}

\begin{verbatim}
## [1] 54.01038
\end{verbatim}

\begin{Shaded}
\begin{Highlighting}[]
\CommentTok{\# round to 2 d.p.}
\NormalTok{mean\_age}\OtherTok{\textless{}{-}}\FunctionTok{round}\NormalTok{(mean\_age, }\DecValTok{2}\NormalTok{)}
\NormalTok{mean\_age}
\end{Highlighting}
\end{Shaded}

\begin{verbatim}
## [1] 54.01
\end{verbatim}
\item
  Median \textbf{(2 marks)}

\begin{Shaded}
\begin{Highlighting}[]
\CommentTok{\# calculate the median age}
\NormalTok{median\_age}\OtherTok{\textless{}{-}}\FunctionTok{median}\NormalTok{(age)}
\NormalTok{median\_age}
\end{Highlighting}
\end{Shaded}

\begin{verbatim}
## [1] 54
\end{verbatim}
\item
  Sample standard deviation \textbf{(2 marks)}

\begin{Shaded}
\begin{Highlighting}[]
\CommentTok{\# calculate the sample standard deviation}
\NormalTok{sd\_age}\OtherTok{\textless{}{-}}\FunctionTok{sd}\NormalTok{(age)}
\NormalTok{sd\_age}
\end{Highlighting}
\end{Shaded}

\begin{verbatim}
## [1] 9.132316
\end{verbatim}

\begin{Shaded}
\begin{Highlighting}[]
\CommentTok{\# round sample standard deviation to 2 d.p.}
\NormalTok{sd\_age}\OtherTok{\textless{}{-}}\FunctionTok{round}\NormalTok{(sd\_age, }\DecValTok{2}\NormalTok{)}
\NormalTok{sd\_age}
\end{Highlighting}
\end{Shaded}

\begin{verbatim}
## [1] 9.13
\end{verbatim}
\end{enumerate}

\begin{center}\rule{0.5\linewidth}{0.5pt}\end{center}

\begin{enumerate}
\def\labelenumi{\arabic{enumi}.}
\setcounter{enumi}{1}
\item
  A researcher wishes to investigate the relationship between gender and
  cholesterol levels.

  Key-points to notice here is that \textbf{we are measuring cholesterol
  levels} based on gender, making \textbf{cholesterol} the
  \emph{dependent variable}, and \textbf{gender} the \emph{independent
  variable}.

  Using the dataset, perform an \emph{appropriate statistical test} in
  \emph{R} and interpret the output. Your answer should include a
  \textbf{null hypothesis}, an \textbf{alternative hypothesis }, and a
  \textbf{p-value } \textbf{(5 marks)}

  \begin{itemize}
  \item
    Since \emph{gender} (categorical) and \emph{cholesterol}
    (continuous) are \textbf{independent} from eachother, we conduct an
    independent (two-sampled) t-test to test the following hypotheses:

    \textbf{Null hypothesis:}

    \begin{itemize}
    \tightlist
    \item
      There is no difference in mean cholesterol levels between genders.
    \end{itemize}

    \textbf{Alternative hypothesis:}

    \begin{itemize}
    \tightlist
    \item
      There is a significant difference in mean cholesterol levels
      between genders.
    \end{itemize}

    \textbf{p-value:}

\begin{Shaded}
\begin{Highlighting}[]
\CommentTok{\# peek at the dataset so we can view the column names}
\FunctionTok{head}\NormalTok{(heart, }\AttributeTok{n=}\DecValTok{2}\NormalTok{)}
\end{Highlighting}
\end{Shaded}

\begin{verbatim}
##   age gender cp trtbps chol fbs restecg maxheartrate exng
## 1  60      1  3    145  233   1       0          150    0
## 2  35      1  2    130  250   0       1          187    0
\end{verbatim}

\begin{Shaded}
\begin{Highlighting}[]
\CommentTok{\# perform a t{-}test on the heart data set between chol and gender}
\FunctionTok{t.test}\NormalTok{(chol}\SpecialCharTok{\textasciitilde{}}\NormalTok{gender, }\AttributeTok{data=}\NormalTok{heart)}
\end{Highlighting}
\end{Shaded}

\begin{verbatim}
## 
##  Welch Two Sample t-test
## 
## data:  chol by gender
## t = 2.9765, df = 133.12, p-value = 0.003464
## alternative hypothesis: true difference in means between group 0 and group 1 is not equal to 0
## 95 percent confidence interval:
##   7.289758 36.169316
## sample estimates:
## mean in group 0 mean in group 1 
##        262.6989        240.9694
\end{verbatim}

    \textbf{- Interpretation:}

    \begin{itemize}
    \item
      The Welch two-sample t-test provides \textbf{evidence to reject
      the null hypothesis}, that \textbf{there will be no difference in
      mean cholesterol levels between genders, \emph{p}=0.003}.
    \item
      This interpretation is further supported with \textbf{a large
      t-value of 2.98.}
    \item
      \textbf{Further testing is recommended} to determine the source or
      magnitude of this cholesterol difference.
    \end{itemize}
  \end{itemize}
\end{enumerate}

\begin{center}\rule{0.5\linewidth}{0.5pt}\end{center}

\begin{enumerate}
\def\labelenumi{\arabic{enumi}.}
\setcounter{enumi}{2}
\item
  A researcher wishes to investigate the predictors of maximum heart
  rate. Using the dataset provided, fit an apprpopriate statistical
  model of your choice. Make sure to include age as one of the
  predictors \textbf{(14 marks)}

  Since \emph{maximum heart rate} and \emph{age} are both continuous
  variables, a regression model will probably be suitable for this
  analysis where we test for a relationship between ones age and their
  maximum heart rate.

  \begin{enumerate}
  \def\labelenumii{\alph{enumii})}
  \item
    What are the null and alternative (two-sided) hypotheses associated
    with age? \textbf{(2 marks)}

    \textbf{Null hypothesis:} H₀: β₁ = 0

    \begin{itemize}
    \item
      The regression coefficient for age is zero (i.e., age has no
      linear effect on maximum heart rate.)

      \subsubsection{Note:}\label{note}

      This means that for every x increase in age, we would not observe
      an increase in y for heart rate.
    \end{itemize}

    \textbf{Alternative hypothesis:} H₁: β₁ ≠ 0

    \begin{itemize}
    \tightlist
    \item
      There will be a significant relationship between age and maximum
      heart-rate.
    \end{itemize}
  \item
    What is the estimated coefficient of age and what does this imply in
    terms of the relationship between age and maximum heart rate?
    \textbf{(3 marks)}
  \item
    Comment on the predictors you have included in the model and justify
    the inclusion of these variables in your model \textbf{(3 marks)}
  \item
    interpret the model output in plain English. Your interpretation
    must include the p-value associated with each coefficient (slopes
    and intercept) test, and R-squared value \textbf{(3 marks)}
  \item
    Comment on whether the assumptions of a linear regression model have
    been met. Include appropriate statistical plots to support your
    answer \textbf{(3 marks)}
  \end{enumerate}
\end{enumerate}

\begin{center}\rule{0.5\linewidth}{0.5pt}\end{center}

\subsection{Data dictionary:}\label{data-dictionary}

cp: chest pain type

\begin{enumerate}
\def\labelenumi{\arabic{enumi}.}
\tightlist
\item
  typical angina
\item
  atypical angina
\item
  non-anginal pain
\item
  asymptomatic
\end{enumerate}

trtbps:

\begin{itemize}
\tightlist
\item
  resting blood pressure (in mm Hg)
\end{itemize}

chol:

\begin{itemize}
\tightlist
\item
  cholesterol (in mg/dl)
\end{itemize}

gender:

\begin{enumerate}
\def\labelenumi{\arabic{enumi}.}
\tightlist
\item
  female;
\item
  male;
\end{enumerate}

fbs: (fasting blood sugar \textgreater{} 120 mg/dl)

\begin{enumerate}
\def\labelenumi{\arabic{enumi}.}
\setcounter{enumi}{-1}
\tightlist
\item
  False
\item
  True
\end{enumerate}

rest\_ecg: resting electrocardiographic result

\begin{enumerate}
\def\labelenumi{\arabic{enumi}.}
\setcounter{enumi}{-1}
\tightlist
\item
  Normal
\item
  Having ST-T wave abnormality (T wave inversions and/or ST elevation or
  depression of \textgreater{} 0.05 mV)
\item
  Showing probably or definite left ventricular hypertrophy by Estes'
  criteria
\end{enumerate}

\begin{center}\rule{0.5\linewidth}{0.5pt}\end{center}

\begin{Shaded}
\begin{Highlighting}[]
\FunctionTok{summary}\NormalTok{(cars)}
\end{Highlighting}
\end{Shaded}

\begin{verbatim}
##      speed           dist       
##  Min.   : 4.0   Min.   :  2.00  
##  1st Qu.:12.0   1st Qu.: 26.00  
##  Median :15.0   Median : 36.00  
##  Mean   :15.4   Mean   : 42.98  
##  3rd Qu.:19.0   3rd Qu.: 56.00  
##  Max.   :25.0   Max.   :120.00
\end{verbatim}

\subsection{Including Plots}\label{including-plots}

You can also embed plots, for example:

\pandocbounded{\includegraphics[keepaspectratio]{A1_files/figure-latex/pressure-1.pdf}}

Note that the \texttt{echo\ =\ FALSE} parameter was added to the code
chunk to prevent printing of the R code that generated the plot.

\end{document}
